\chapter{Schwarz Inequalities and Multiplicative Domains: Supporting Results}
\label{ch:schwarz_supporting_results}

This appendix supports Chapter~\ref{ch:schwarz}. It records the rank-one
rigidity argument used to identify the scalar factor in a positive retraction.

The following theorem supplies the scalarity step in
Theorem~\ref{thm:positive_retraction_one_matrix_factor}.

\begin{theorem}[Rank-one ray scalar rigidity]
    \label{thm:rankone_ray_scalar_rigidity}
    \lean{WolfProps.exists_eq_smul_id_of_maps_rankOne_to_span}
    \leanok
    Let $T:\MN{D}\to\MN{D}$ be complex-linear. Suppose that for every
    $v\in\C^D$ there is a scalar $c_v\in\C$ such that
    $T(\ketbra{v}{v})=c_v\ketbra{v}{v}$.
    Then there is a scalar $c\in\C$ such that $T=c\,\id$.
\end{theorem}

\begin{proof}
    \leanok
    Write $P_v=\ketbra{v}{v}$. If $D=0$, then $\MN{D}$ is the zero vector
    space, so the conclusion holds with $c=0$. If $D=1$, the single matrix unit
    is $P_{e_0}$. The hypothesis gives $T(P_{e_0})=c_{e_0}P_{e_0}$, and
    linearity gives $T=c_{e_0}\,\id$ on all of $\MN{1}$.

    It remains to consider $D\geq2$. Set $c=c_{e_0}$. For $i\ne j$, the identities
    \begin{align}
        P_{e_i+e_j}+P_{e_i-e_j}
         & =2P_{e_i}+2P_{e_j},
        \notag                              \\
        c_{e_i+e_j}P_{e_i+e_j}+c_{e_i-e_j}P_{e_i-e_j}
         & =2c_{e_i}P_{e_i}+2c_{e_j}P_{e_j}
        \notag
    \end{align}
    give, from their $(i,i)$, $(j,j)$, and $(i,j)$ entries,
    $c_{e_i+e_j}=c_{e_i-e_j}=c_{e_i}=c_{e_j}$.
    The imaginary parallelogram identity and its image under $T$ are
    \begin{align}
        P_{e_i+\mathrm{i}e_j}+P_{e_i-\mathrm{i}e_j}
         & =2P_{e_i}+2P_{e_j},
        \notag                               \\
        c_{e_i+\mathrm{i}e_j}P_{e_i+\mathrm{i}e_j}
        +c_{e_i-\mathrm{i}e_j}P_{e_i-\mathrm{i}e_j}
         & =2c_{e_i}P_{e_i}+2c_{e_j}P_{e_j}.
        \notag
    \end{align}
    Comparing the same three entries gives
    $c_{e_i+\mathrm{i}e_j}=c_{e_i-\mathrm{i}e_j}=c_{e_i}=c_{e_j}$.
    Taking the pair $(0,i)$ shows that $c_{e_i}=c$ for every $i$.
    The hypothesis therefore gives $T(E_{ii})=cE_{ii}$ for the diagonal matrix
    units. For $i\ne j$, the polarization identity
    \begin{align}
        4E_{ij}
         & =P_{e_i+e_j}-P_{e_i-e_j}
        +\mathrm{i}P_{e_i+\mathrm{i}e_j}
        -\mathrm{i}P_{e_i-\mathrm{i}e_j}
        \notag
    \end{align}
    gives $T(E_{ij})=cE_{ij}$. Thus $T=c\,\id$ on every matrix unit.
\end{proof}

Thus rank-one domination in
Theorem~\ref{thm:positive_retraction_one_matrix_factor} produces one scalar
functional rather than a vector-dependent family of coefficients.
